\chapter{Fejlesztői dokumentáció}\label{ch:architecture}
\section{Tervezés}\label{sec:design}

\subsection{Feladat részletes leírása}

Az alkalmazás kettő fő részből áll: egy ASP.NET Core alapú backend szolgáltatásból, egy Vue és TypeScript alapú frontend egyoldalas alkalmazásból. Fontos hogy ezen részek önmagukban is megállják a helyüket, de együtt egy teljes funkcionalitást nyújtó rendszert alkotnak. Ennek megvalósásában az stabil API-kapcsolat, a jól definiált adatmodell, valamint a felhasználói élményre fókuszáló frontend kialakítása a fő célok közé tartozik. 

\subsection{Funkcionális követelmények}\label{ch:requirements}

Funcionális követelmények tekintetében a projekt fő területei a következők:

\begin{itemize}
	\item Felhasználókezelés: regisztráció, bejelentkezés, profil megjelenítése és módosítása.
	\item Receptkezelés: új recept létrehozása, szerkesztése, törlése, részletes receptnézet.
	\item Keresés és szűrés: receptlista böngészése, kulcsszavas keresés, szűrés allergének alapján.
	\item Közösségi funkciók: receptek értékelése, kommentek rögzítése és törlése.
\end{itemize}

\subsection{Nem funkcionális követelmények}

A rendszerrel szemben a nem funkcionális követelmények a következők:

\begin{itemize}
	\item átlátható és reszponzív felhasználói felület,
	\item jól dokumentált, stabilan használható HTTP API,
	\item olyan adatmodell, amely a recept-, komment- és felhasználói adatok közötti kapcsolatokat konzisztensen kezeli,
	\item könnyen továbbfejleszthető frontend és  backend oldalon,
	\item fejlesztői munkát segítő konténeres indítási és build környezet \cite{docker}.
	\item a rendszer helyi környezetben való könnyű indíthatósága, minimális külső függőséggel.
\end{itemize}

\subsection{Felhasználói esetek}

% TODO: Use case diagram és leírások, a követelmények és a megvalósított funkciók tételes összerendelése
Az alábbi ábrán jól látható a rendszer főbb funkciói és azok kapcsolatai. A diagramon a főbb szereplők (felhasználó, adminisztrátor) és azok interakciói jelennek meg a rendszer különböző részeivel. Minden adminisztrátor egyben az alap felhasználó funkcióival is rendelkezik, ez külön nincs feltüntetve a diagramon a könnyebb átláthatóság érdekében.

\begin{figure}[H]
	\centering
	\includegraphics[width=\textwidth]{usecases.png}
	\caption{Felhasználói esetek diagramja}\label{fig:usecase-diagram}
\end{figure}

A felhasználói történetek a~\ref{ch:stories} függelékben találhatóak meg. Részletesen, táblázatban összefoglalva.

\subsection{Architekturális áttekintés}

A projekt réteges, webalkalmazás-jellegű architektúrára épül. A felhasználó egy böngészőben futó kliensalkalmazással lép kapcsolatba, amely HTTP hívásokon keresztül éri el a backend szolgáltatást. A backend végzi az üzleti logika, az adatelérés, a hitelesítéshez kapcsolódó folyamatok, valamint az API-dokumentáció kiszolgálásának fő részét. Az adatok perzisztálását egy relációs adatbázis biztosítja~\cite{vue,aspnetcore,postgresql}.



A fő komponensek:

\begin{itemize}
	\item Frontend: Vue 3 alapú egyoldalas alkalmazás, amely a felhasználói folyamatokat és az oldalak közti navigációt kezeli.
	\item API kliens: OpenAPI alapján generálható klienscsomag, amely a frontend és a backend közötti szerződést teszi követhetővé.
	\item Backend: ASP.NET Core Web API, amely kontroller-, szolgáltatás- és repository-szintű felelősségekre tagolódik
	\item Adatbázis: PostgreSQL  adatbázis, amely a receptek, kommentek, felhasználói adatok és egyéb kapcsolódó információk tárolásáért felelős.
	\item Proxy: Nginx alapú reverse proxy, amely a frontend és a backend szolgáltatásokat egyetlen végpont alá helyezi
	\item Adminer: könnyű adatbázis-kezelő eszköz, adatbázis tartalmának gyors áttekintését és módosítását teszi lehetővé
	\item Docker volumes: az adatbázis adatai és a képeket egy-egy megosztott kötetben tároljuk.
\end{itemize}


\subsection{REST API tervezés}

Az alkalmazás tervezésénél fontos szempont volt, hogy a backend és a frontend önmagukban is teljes értékű, jól működő komponensek legyenek. Ennek eléréséhez elengedhetetlen, hogy a két rendszer jól dokumentált és következetes kommunikációt folytasson egymással. Ez lehetővé teszi a frontend és a backend párhuzamos fejlesztését, valamint azt, hogy a komponensek külön-külön is könnyen tesztelhetők legyenek.

Erre a problémára a Representational State Transfer (REST) API tervezési megközelítés nyújt megoldást, amely egy széles körben alkalmazott szabvány a webes szolgáltatások kialakításában. A REST API-k használata biztosítja a frontend és a backend közötti egyszerű, egységes és hatékony kommunikációt.

A megközelítés alkalmazása számos előnnyel jár: támogatja a rendszer skálázhatóságát, megkönnyíti a tesztelést, valamint lehetővé teszi a komponensek egymástól független fejlesztését és karbantartását. A backend szolgáltatás nem függ közvetlenül a frontend implementációjától, így a két réteg lazán csatolt módon működik együtt.

További előny, hogy a rendszer könnyen kiterjeszthető más platformokra is. Amennyiben a jövőben mobilalkalmazás vagy egyéb kliens kerül bevezetésre, ezek ugyanazt az API-t használhatják, így elegendő csupán a megjelenítési réteget az adott platformhoz igazítani.

Előnyök mellett sajnos hátrányokkal is jár. REST API-k tervezése időigényes lehet, különösen összetett üzleti logika esetén. Emellett növeli komplexitást, több eszközt kell használni fejlesztéshez és több programozási nyelvet kell ismerni.

Az api leírására az OpenAPI specifikációt választottam. Ez egy széles körben használt szabvány, amely lehetővé teszi a REST API-k részletes dokumentálását. A HTTP végpontok leírása a követekezőképpen néz ki:

\begin{figure}[H]
	\centering
	\includegraphics[width=\textwidth]{recipes-endpoints.png}
	\caption{Receptek HTTP végpontjai}\label{fig:recipe-endpoints-diagram}
\end{figure}

\begin{figure}[H]
	\centering
	\includegraphics[width=\textwidth]{ingredients-endpoints.png}
	\caption{Hozzávalók HTTP végpontjai}\label{fig:ingredients-endpoints-diagram}
\end{figure}

\begin{figure}[H]
	\centering
	\includegraphics[width=\textwidth]{profile-endpoints.png}
	\caption{Felhasználói profil HTTP végpontjai}\label{fig:profile-endpoints-diagram}
\end{figure}

\begin{figure}[H]
	\centering
	\includegraphics[width=\textwidth]{auth-endpoints.png}
	\caption{Hitelesítési HTTP végpontok}\label{fig:auth-endpoints-diagram}
\end{figure}

\subsection{Adatbázisterv}

% TODO: ER diagram, fő entitások és kapcsolataik, Entity Framework Core migrációk

\subsection{Felhasználói felület terve}

% TODO: Wireframe-ek, a felhasználói felület terveinek bemutatása, a fő oldalak és folyamatok áttekintése